\section{Design}
\label{sec:design}

Figure~\ref{fig:overview} shows the workflow. The sources on the left are what
git tracks; everything on the right is produced by one \texttt{paperkit build}
and stays out of version control.

\begin{figure}[t]
  \centering
  % The workflow figure, as TikZ: typeset with the paper's fonts, versioned as
% text, and nothing to convert. Sized to the column: the report counts an
% overfull line, and a figure a few points too wide is one.
\begin{tikzpicture}[
    node distance=3mm and 5mm,
    box/.style={draw, rounded corners=1pt, minimum width=20mm, minimum height=6mm,
      font=\scriptsize, align=center},
    arr/.style={-{Stealth[length=2mm]}, thick},
  ]
  \node[box] (src) {sources\\paper.tex, content/,\\header/, figures/};
  \node[box, right=of src] (build) {paperkit build};
  \node[box, above right=of build] (pdf) {PDF, log};
  \node[box, right=of build] (report) {report};
  \node[box, below right=of build] (txt) {title.txt,\\abstract.txt};
  \node[box, below=of build] (zip) {paperkit zip};
  \node[box, right=of zip] (archive) {paper-sources.zip};
  \draw[arr] (src) -- (build);
  \draw[arr] (build) -- (pdf);
  \draw[arr] (build) -- (report);
  \draw[arr] (build) -- (txt);
  \draw[arr] (src) |- (zip);
  \draw[arr] (zip) -- (archive);
\end{tikzpicture}

  % acmart requires a \Description for every figure, for screen readers, and
  % warns about each one that lacks it.
  \Description{Sources on the left feed paperkit build, which produces the
  PDF, the report and the text exports; paperkit zip produces the archive.}
  \caption{One build produces the PDF, the build log, the exported title and
  abstract, and the report; one \texttt{paperkit zip} produces the archive.
  Nothing on the right is committed.}
  \label{fig:overview}
\end{figure}

\paragraph{Sources in directories}
The main file \texttt{paper.tex} names the pieces and their order; the text
lives in \texttt{content/}, numbered so that the directory listing reads like
the table of contents, and the preamble in \texttt{header/}.
Figure~\ref{fig:overview} is itself a \TeX\ source, kept with the other
figures, so it is typeset with the paper's fonts and needs no conversion.

\paragraph{Two targets}
The supplementary material is a second main file in the same tree. It shares
\texttt{header/} and the bibliography, so a macro renamed in one place is
renamed in both, and an entry cited only in the appendix is not reported as
unused when the paper is checked.

\paragraph{A pinned environment}
The configuration names a dated \TeX\ Live snapshot and its checksum, so the
image that builds the paper has the same packages on every rebuild, on every
machine, and the build timestamp comes from the commit, so two builds of one
commit produce the same bytes~\cite{reproducible}.
